% !TEX root = ../main.tex

\section{Resultados}
\label{sec:results}

\subsection{Protocolo de validaci\'on}
\label{sec:results:validation}

La validaci\'on empirica se realiz\'o en seis sesiones con
caracterizados en distintos perfiles: novatos sin
experiencia previa en programaci\'on, programadores expertos con
a\~nos de experiencia, estudiantes de secundaria, ni\~nos y
profesionales del campo art\'istico. La diversidad de perfiles fue
intencional. \sculpt aspira a ser inclusivo, y la validaci\'on debe
ponerlo a prueba con pre\'ambulos cognitivos y motivacionales muy
distintos. El objetivo de estas sesiones de validación, con respecto
a los perfiles de los usuarios es: (1) mostrar la posibilidad de \sculpt
como un lenguaje de programación que rompe las barreras de accesibilidad
para nuevos usuarios. (2) La efectividad de \sculpt como lenguaje de programación,
permitiendo a programadores con experiencia construir programas completos.
(3) El uso de la programación como un medio de expresión artistico, donde
artistas pueden construir esculturas que son a su vez programas computables. 

Las sesiones contaron con 1 a 3 facilitadores cuya labor consisti\'o en
guiar las presentaciones, resolver dudas y supervisar el progreso.
Cada sesi\'on sigui\'o cuatro fases. (1) Una \emph{introducci\'on}
de 30 a 45~minutos dirigida por los facilitadores, adaptada al perfil del grupo: m\'as detallada
sobre los componentes f\'isicos para los novatos y m\'as t\'ecnica
para los expertos. (2) Una fase de \emph{exploraci\'on libre} de
una hora, sin tareas asignadas, en la que los participantes
manipulan las piezas a su ritmo para familiarizarse con sus formas
y conexiones (3) Una fase de \emph{tareas computacionales} de
dos horas, con conjuntos de problemas calibrados al nivel de cada
grupo (\eg sumas iteradas, condicionales para romper bucles); los
facilitadores resolv\'ian dudas pero el participante completaba la
tarea por su cuenta. (4) Una \emph{encuesta y entrevista} final opcional
sobre usabilidad, accesibilidad y experiencia general; en el caso
de menores de edad, la encuesta tambi\'en la respondieron los
docentes acompa\~nantes, con preguntas adicionales sobre la
viabilidad de \sculpt como herramienta pedag\'ogica.

El an\'alisis combin\'o m\'etodos cualitativos y cuantitativos.
En lo cuantitativo se midi\'o el tiempo a tarea y el n\'umero de
errores cometidos durante la fase de programaci\'on. En lo
cualitativo se analizaron las entrevistas y las preguntas abiertas
de la encuesta, buscando patrones recurrentes en la percepci\'on
de las piezas, en los obst\'aculos sint\'acticos y en la
experiencia est\'etica de manipular el lenguaje. Adicionalmente,
con el subgrupo de artistas se realizaron sesiones extra
dedicadas exclusivamente a la creaci\'on de piezas no funcionales,
con el fin de evaluar a \sculpt como medio expresivo independiente
de su capacidad computacional.

\subsection{Validaci\'on emp\'irica}
\label{sec:results:empiricalvalidation}

\subsubsection{Usabilidad}
La fisicalidad de las piezas fue el factor m\'as recurrente en los
reportes positivos, especialmente en los grupos j\'ovenes, que
identificaron y diferenciaron los componentes con facilidad. La
integridad material de las piezas fue un punto de atenci\'on: el PLA es
un pl\'astico relativamente blando, propenso a deformaci\'on con uso
repetido. Algunas piezas se perdieron a lo largo de las sesiones y
otras mostraron desgaste, pero el conjunto resisti\'o el ciclo de uso
sin afectar la continuidad de las pruebas.

\subsubsection{Programaci\'on}
Aproximadamente el 50\,\% del total de participantes complet\'o las
tareas computacionales asignadas. El desempe\~no por sub-poblaci\'on,
sin embargo, no es el esperable. La Tabla~\ref{tab:validation} resume los
hallazgos por grupo.

\begin{table}[h]
\centering
\small
\caption{Hallazgos de las tareas computacionales por sub-poblaci\'on.}
\label{tab:validation}
\begin{tabular}{@{}lll@{}}
\toprule
Sub-poblaci\'on & Completitud & Observaci\'on cualitativa \\
\midrule
Ni\~nos                 & Baja      & Excelente manejo f\'isico; dificultad conceptual. \\
Adolescentes novatos    & Media     & Completaron en $\sim$1{,}5\,h promedio. \\
Adultos novatos         & Baja      & Comprensi\'on del lenguaje sin completar tareas. \\
Programadores expertos  & Baja      & Dificultad similar a novatos. \\
Artistas                & Nula      & Comprenden las piezas, no completan tareas. \\
\bottomrule
\end{tabular}
\end{table}

El resultado m\'as inesperado fue la dificultad de los programadores
expertos: pese a su experiencia previa, no superaron a los
adolescentes novatos en completitud. Una lectura plausible es que el
modelo mental que aporta la experiencia con lenguajes textuales no
transfiere directamente al razonamiento sobre pilas LIFO ensambladas
f\'isicamente. Los adolescentes, sin esos modelos previos, abordaron
el lenguaje en sus propios t\'erminos.

\subsubsection{Datos de la encuesta}
La encuesta cuantitativa (escala Likert 1--5) se aplic\'o a un
sub-conjunto de $n=25$ participantes de forma voluntaria para los que se obtuvieron
respuestas completas. 
Las preguntas se agruparon en cinco bloques: (1) informaci\'on
demogr\'afica y de experiencia previa, (2) aspectos cuantitativos del
lenguaje (\eg dificultad, claridad), (3) aspectos cualitativos del lenguaje
(\eg disfrute, creatividad, percepcion de utilidad), (4) evaluaci\'on comparativa de conocimientos
de prigamaci\'on, y (5) percepci\'on de docentes/evaluadores sobre el potencial pedag\'ogico de \sculpt.
La Tabla~\ref{tab:survey-likert} resume los rangos
de media por bloque de preguntas.

\begin{table}[h]
\centering
\small
\caption{Agregados de la encuesta cuantitativa ($n=25$, escala 1--5).}
\label{tab:survey-likert}
\begin{tabular}{@{}lc@{}}
\toprule
Bloque de preguntas & Media (rango) \\
\midrule
Aspectos cuantitativos del lenguaje              & 2{,}2 -- 3{,}4 \\
Evaluaci\'on comparativa frente a lenguajes textuales & 2{,}8 -- 3{,}5 \\
Percepci\'on de docentes/evaluadores             & 4{,}0 -- 4{,}3 \\
\bottomrule
\end{tabular}
\end{table}

La diferencia entre los bloques es informativa. La percepci\'on
t\'ecnica de los participantes es media-baja, consistente con la
dificultad reportada para completar las tareas. La percepci\'on de
los docentes que acompa\~naron las sesiones es claramente positiva,
lo que sugiere viabilidad pedag\'ogica incluso cuando el desempe\~no
inmediato de los aprendices es modesto. Las respuestas cualitativas
refuerzan ambos lados: los participantes destacan que ``llevar
programaci\'on al mundo 3D es una haza\~na'' y que el lenguaje es
``interesante y creativo para ni\~nos peque\~nos'', pero tambi\'en
reportan que ``los bloques a veces generan confusi\'on'' y que
``los conceptos de operaciones son confusos''. Las sugerencias
m\'as frecuentes apuntan a material did\'actico de apoyo
(diapositivas, ejercicios incrementales) y a integraci\'on con
dominios concretos como videojuegos o matem\'aticas.

\subsubsection{Est\'etica}
Las sesiones con artistas dieron los resultados m\'as positivos en
percepci\'on. Los participantes valoraron la manipulaci\'on tangible
y reportaron que ``la fisicalidad fue quintaesencial a su
experiencia''. Aunque no completaron las tareas asignadas,
expresaron satisfacci\'on con el proceso e inter\'es por seguir
usando el lenguaje, al punto de solicitar sesiones adicionales en
las que construyeron piezas no funcionales con prop\'osito
puramente est\'etico. El uso emergente que dieron a las piezas, como
ensambles no previstos en la especificaci\'on, exploraciones
de la forma por la forma, y colaboraciones entre escultura e ilustración valida la apertura est\'etica como
caracter\'istica diferenciadora del lenguaje.

\endinput
